\section{Feature implementation}

\subsection{Destinations and quests}
Two versioned launch packs ship in application assets: Nha Trang and Ho Chi Minh
City. Each pack contains city metadata, route plans, public landmarks, photo
quests, fallback instructions, points of interest, and review metadata. The
repository validates identifiers and relationships while decoding the JSON, so a
malformed pack becomes an understandable loading failure rather than corrupting
a run.

Quest eligibility is evaluated from the latest accepted location and the quest's
radius. When GPS is unavailable or unreliable, the content-defined fallback may
allow the user to confirm physical presence explicitly. A required caption is
validated before evidence is committed. This balances location awareness with
the practical limits of urban GPS and device testing.

\appscreenshot{images/route-details.png}{105mm}{Route details present metrics, safety notes, landmarks, and photo quests before tracking begins.}

\subsection{Run lifecycle and location processing}
The session lifecycle is represented explicitly as idle, active, paused, and
finished states. Only active intervals contribute to duration. Pausing closes the
current segment, while resuming opens a new segment. Finish derives a summary and
clears the active checkpoint only after the final data has been saved.

Consecutive accepted points use the Haversine formula:
\[
a=\sin^2(\Delta\varphi/2)+\cos(\varphi_1)\cos(\varphi_2)
  \sin^2(\Delta\lambda/2),\qquad
d=2R\arctan2(\sqrt a,\sqrt{1-a}).
\]
Implausible jumps are rejected using distance, elapsed time, and calculated
speed. Non-increasing timestamps are invalid. Pace is unavailable for zero
distance rather than represented by an infinite or misleading number. Calories
are clearly presented as an estimate.

The Android location adapter requests high accuracy, applies a five-meter
distance filter, and configures a foreground notification for an ongoing run.
The manifest declares fine/coarse location and the location foreground-service
type. Tracking still begins only after a user-visible action.

\appscreenshot{images/tracker.png}{108mm}{The active tracker keeps lifecycle controls, live metrics, and quest actions visible.}

\subsection{Camera evidence and retained media}
Camera capture paths are temporary. After a valid capture and caption, the photo
store copies the image into application-controlled support storage and the
evidence record refers to that retained path. Cancellation, permission failure,
and plugin errors return to a usable screen. History deletion removes dependent
story documents before the run record and cleans retained evidence only after
repository records are consistent.

\subsection{History and achievements}
Journey combines completed runs, editable story projects, and personal
achievements. Empty, loading, populated, and failure states are explicit. Run
cards reopen summaries; story cards reopen their source project; destructive
deletion requires confirmation. Achievement criteria are deterministic and
personal, avoiding a public leaderboard or account requirement.

\appscreenshot{images/history.png}{105mm}{Journey preserves completed runs, editable story projects, and personal achievements.}

\subsection{Story Studio and export}
Story Studio stores an ordered document with a fixed 1080 by 1920 canvas. Original
templates provide starting compositions but remain ordinary editable documents.
Elements carry their own type, content, style, layer order, and transform.
Selection, move, scale, rotate, duplicate, delete, bring-forward, send-backward,
undo, and redo operations produce new immutable documents.

Export renders from the saved document model rather than taking a screenshot of
the editor. Selection borders, handles, and guides are not part of the rendered
canvas. The output dimensions are checked before the PNG is offered to the
Android share sheet. A failed renderer or share service produces recoverable
feedback and does not destroy the project.

\appscreenshot{images/story-studio.png}{110mm}{Story Studio turns run evidence into a layered, editable portrait composition.}
